\usepackage{tikz}
\usepackage{pgfplots}
\usepgfplotslibrary{fillbetween}
\usetikzlibrary{patterns, patterns.meta, positioning, math}

\pgfplotsset{compat=1.18}

\usepackage{tikz-3dplot}
\usetikzlibrary{3d,calc}

\pgfplotsset{
	twoDaxis/.style = {
			no marks,
			samples = 100,
			axis lines = center,
			xlabel = {$x$},
			ylabel = {$y$},
			every axis x label/.style={
					at={(ticklabel* cs:1.05)},
					anchor=west,
				},
			every axis y label/.style={
					at={(ticklabel* cs:1.05)},
					anchor=south,
				},
		},
	threeDaxis/.style = {
			axis z line = none,
			axis lines=center, % remove this and you get weird box behavior, not good for calc 2 students
			xlabel=$x$, % you know this one
			every axis x label/.style={
					at={(ticklabel* cs:1.05)},
					anchor=west,
				}, % place the xlabel in a better spot
			ylabel=$y$, % you know this one
			every axis y label/.style={
					at={(ticklabel* cs:1.05)},
					anchor=north,
				}, % places y label in a better spot
		},
	threeDaxiswithyzswap/.style = {
			axis y line = none,
			axis lines=center, % remove this and you get weird box behavior, not good for calc 2 students
			xlabel=$x$, % you know this one
			every axis x label/.style={
					at={(ticklabel* cs:1.05)},
					anchor=west,
				}, % place the xlabel in a better spot
			zlabel=$y$, % you know this one
			every axis z label/.style={
					at={(ticklabel* cs:1.05)},
					anchor=south, % which side of the label gets pinned to the point 5% past the axis end point?
				}, % places y label in a better spot
			xticklabel style = {font=\tiny, yshift=1ex},
			zticklabel style = {font=\tiny, xshift=1ex},
		}
}
% axis direction cs makes it only work in pgfplots
\def\centerarc[#1](#2)(#3:#4:#5)% Syntax: [draw options] (center) (initial angle:final angle:radius)
{ \draw[#1] ($(#2)+(axis direction cs:{#5*cos(#3)},{#5*sin(#3)})$) arc (#3:#4:#5); }
\def\centerarc[#1](#2)(#3:#4:#5:#6)% Syntax: [draw options] (center) (initial angle:final angle:xradius:yradius)
{ \draw[#1] ($(#2)+(axis direction cs:{#5*cos(#3)},{#6*sin(#3)})$) arc  [start angle = #3, end angle = #4, x radius = #5, y radius = #6]; }
\def\diskslicevert[#1](#2)(#3:#4)% Syntax: [draw options] (center) (xradius:yradius) use {} if functions on things!!!!!!
{ 		\draw [<->, #1] (#2) -- ($(#2)+(axis direction cs:{0},{#4})$) node [midway, right] {$R$};
	\centerarc[#1](#2)(20:340:{#3}:{#4});
	\centerarc[dotted, #1](#2)(340:380:{#3}:{#4});
}
\def\diskslicevert[#1,r=#2](#3)(#4:#5)% Syntax: [draw options] (center) (xradius:yradius)
{ 		\draw [<->, #1] (#3) -- ($(#3)+(axis direction cs:{0},{#5})$) node [midway, right] {#2};
	\centerarc[#1](#3)(20:340:{#4}:{#5});
	\centerarc[dotted, #1](#3)(340:380:{#4}:{#5});
}
\def\diskslicevertflip[#1,r=#2](#3)(#4:#5)% Syntax: [draw options] (center) (xradius:yradius)
{ 		\draw [<->, #1] (#3) -- ($(#3)+(axis direction cs:{0},{-#5})$) node [midway, right] {#2};
	\centerarc[#1](#3)(20:340:{#4}:{#5});
	\centerarc[dotted, #1](#3)(340:380:{#4}:{#5});
}
\def\diskslicehor[#1](#2)(#3:#4)% Syntax: [draw options] (center) (xradius:yradius)
{ 	\draw [<->, #1] (#2) -- ($(#2)+(axis direction cs:{#3},{0})$) node [midway, above] {$R$};
	\centerarc[#1](#2)(0:70:#3:#4);
	\centerarc[dotted, #1](#2)(70:110:#3:#4);
	\centerarc[#1](#2)(110:360:#3:#4);
}
\def\diskslicehor[#1,r=#2](#3)(#4:#5)% Syntax: [draw options] (center) (xradius:yradius)
{ 	\draw [<->, #1] (#3) -- ($(#3)+(axis direction cs:{#4},{0})$) node [midway, above] {#2};
	\centerarc[#1](#3)(0:70:#4:#5);
	\centerarc[dotted, #1](#3)(70:110:#4:#5);
	\centerarc[#1](#3)(110:360:#4:#5);
}
\def\diskslicehorflip[#1,r=#2](#3)(#4:#5)% Syntax: [draw options] (center) (xradius:yradius)
{ 	\draw [<->, #1] (#3) -- ($(#3)+(axis direction cs:{-#4},{0})$) node [midway, above] {#2};
	\centerarc[#1](#3)(0:70:#4:#5);
	\centerarc[dotted, #1](#3)(70:110:#4:#5);
	\centerarc[#1](#3)(110:360:#4:#5);
}
\def\boxfill[#1](#2,#3,#4)(#5,#6,#7){ % suggested #1 [fill opacity=0.4,	draw=black,	fill=blue,]
	% top
	\fill[#1] (#2,#3,#7) -- (#2,#6,#7) -- (#5,#6,#7) -- (#5,#3,#7) -- cycle;
	% bottom
	\fill[#1] (#2,#3,#4) -- (#2,#6,#4) -- (#5,#6,#4) -- (#5,#3,#4) -- cycle;
	% front
	\fill[#1] (#2,#3,#7) -- (#2,#3,#4) -- (#5,#3,#4) -- (#5,#3,#7) -- cycle;
	% back
	\fill[#1] (#2,#6,#7) -- (#2,#6,#4) -- (#5,#6,#4) -- (#5,#6,#7) -- cycle;
	% right
	\fill[#1] (#5,#3,#7) -- (#5,#6,#7) -- (#5,#6,#4) -- (#5,#3,#4) -- cycle;
	% left
	\fill[#1] (#2,#3,#7) -- (#2,#6,#7) -- (#2,#6,#4) -- (#2,#3,#4) -- cycle;
}